\documentclass[reqno]{amsart}
\usepackage[T1]{fontenc}
\usepackage{amsmath,amssymb,amsthm}
\usepackage{xr-hyper}
\usepackage[hidelinks]{hyperref}
\externaldocument[hn-]{20261125T013500--hecke-nondegeneracy-on-the-tree__workbook}
\externaldocument[di-]{20261125T020100--divisibility-at-inert-primes-in-general__workbook}
\externaldocument[ea-]{20261126T020130--nondegeneracy-at-an-edge__workbook}

\newcounter{chap}
% Chapter numbers, generated from 20261125T012000--grh-workbook__guide.org
% by scripts/make-chapters.py.  Do not edit: change the order of the
% index and rerun the script (./preview-site.sh and CI do this).
%
% Usage in a sheet, after \newcounter{chap}:
%   % Chapter numbers, generated from 20261125T012000--grh-workbook__guide.org
% by scripts/make-chapters.py.  Do not edit: change the order of the
% index and rerun the script (./preview-site.sh and CI do this).
%
% Usage in a sheet, after \newcounter{chap}:
%   % Chapter numbers, generated from 20261125T012000--grh-workbook__guide.org
% by scripts/make-chapters.py.  Do not edit: change the order of the
% index and rerun the script (./preview-site.sh and CI do this).
%
% Usage in a sheet, after \newcounter{chap}:
%   \input{chapters}\chapnumber{<sheet ID>}
\makeatletter
\providecommand\chapters@set[2]{\@namedef{chapters@#1}{#2}}
\chapters@set{20261125T001000}{1}% zeros-to-ramp-sums
\chapters@set{20261125T001100}{2}% inert-primes-and-estimate-b
\chapters@set{20261125T005000}{3}% galois-fourier-and-averaging
\chapters@set{20261125T005500}{4}% columns-from-a-sphere
\chapters@set{20261125T000000}{5}% slope-method
\chapters@set{20261126T030100}{6}% a-twist-for-each-prime
\chapters@set{20261125T003000}{7}% the-chain-and-the-square-root
\chapters@set{20261125T010000}{8}% rank-lemma-and-zero-estimate
\chapters@set{20261125T010500}{9}% poincare-series-and-near-orthogonality
\chapters@set{20261124T223000}{10}% elliptic-curves-for-the-proof
\chapters@set{20261125T009500}{11}% fock-model-and-theta-intertwiners
\chapters@set{20261124T223100}{12}% complex-multiplication-and-inert-reduction
\chapters@set{20261125T008000}{13}% one-cm-determinant
\chapters@set{20261124T214500}{14}% supersingular-theta-divisibility
\chapters@set{20261124T214600}{15}% the-isotropic-plane
\chapters@set{20261125T013700}{16}% pairings-on-a-formal-group-of-height-two
\chapters@set{20261125T013500}{17}% hecke-nondegeneracy-on-the-tree
\chapters@set{20261126T030200}{18}% the-bound-is-attained
\chapters@set{20261125T009000}{19}% kronecker-limit-and-chowla-selberg
\chapters@set{20261125T002000}{20}% definite-theta-lifts
\chapters@set{20261125T015500}{21}% siegel-weil-in-the-smallest-case
\chapters@set{20261125T002100}{22}% definite-theta-lifts-and-norms
\chapters@set{20261125T016500}{23}% hecke-operators-and-theta-lifts
\chapters@set{20261125T016700}{24}% divisibility-along-the-chain
\chapters@set{20261125T004000}{25}% denominators-and-sha
\chapters@set{20261125T004100}{26}% tate-twists-and-sha
\chapters@set{20261126T020430}{27}% the-minus-part-by-gauss-sums
\chapters@set{20261125T014100}{28}% congruence-modules
\chapters@set{20261125T014500}{29}% eisenstein-congruences-and-selmer-classes
\chapters@set{20261125T014700}{30}% pseudocharacters-and-the-length-bound
\chapters@set{20261125T011000}{31}% galois-representations-black-box
\chapters@set{20261126T030300}{32}% the-data
\chapters@set{20261126T020210}{33}% bessel-poincare-series-and-their-gram-operators
\chapters@set{20261126T030400}{34}% the-seesaw-identity
\chapters@set{20261126T030500}{35}% divisibility-of-the-composite-lift
\chapters@set{20261126T030600}{36}% theorem-a-short-forms
\chapters@set{20261126T030700}{37}% the-determinant-of-a-chain
\chapters@set{20261126T030800}{38}% the-cost-of-one-lift
\chapters@set{20261126T030900}{39}% denominators-of-projections-to-lifts
\chapters@set{20261126T031000}{40}% twists-with-few-congruences
\chapters@set{20261126T031100}{41}% theorem-c-and-the-main-inequality
\providecommand\chapnumber[1]{%
  \@ifundefined{chapters@#1}%
    {\PackageError{chapters}{Sheet #1 is not listed in the index}%
      {List it in 20261125T012000--grh-workbook__guide.org\MessageBreak
       and run scripts/make-chapters.py.}}%
    {\edef\thechap{\@nameuse{chapters@#1}}}}
\makeatother
\chapnumber{<sheet ID>}
\makeatletter
\providecommand\chapters@set[2]{\@namedef{chapters@#1}{#2}}
\chapters@set{20261125T001000}{1}% zeros-to-ramp-sums
\chapters@set{20261125T001100}{2}% inert-primes-and-estimate-b
\chapters@set{20261125T005000}{3}% galois-fourier-and-averaging
\chapters@set{20261125T005500}{4}% columns-from-a-sphere
\chapters@set{20261125T000000}{5}% slope-method
\chapters@set{20261126T030100}{6}% a-twist-for-each-prime
\chapters@set{20261125T003000}{7}% the-chain-and-the-square-root
\chapters@set{20261125T010000}{8}% rank-lemma-and-zero-estimate
\chapters@set{20261125T010500}{9}% poincare-series-and-near-orthogonality
\chapters@set{20261124T223000}{10}% elliptic-curves-for-the-proof
\chapters@set{20261125T009500}{11}% fock-model-and-theta-intertwiners
\chapters@set{20261124T223100}{12}% complex-multiplication-and-inert-reduction
\chapters@set{20261125T008000}{13}% one-cm-determinant
\chapters@set{20261124T214500}{14}% supersingular-theta-divisibility
\chapters@set{20261124T214600}{15}% the-isotropic-plane
\chapters@set{20261125T013700}{16}% pairings-on-a-formal-group-of-height-two
\chapters@set{20261125T013500}{17}% hecke-nondegeneracy-on-the-tree
\chapters@set{20261126T030200}{18}% the-bound-is-attained
\chapters@set{20261125T009000}{19}% kronecker-limit-and-chowla-selberg
\chapters@set{20261125T002000}{20}% definite-theta-lifts
\chapters@set{20261125T015500}{21}% siegel-weil-in-the-smallest-case
\chapters@set{20261125T002100}{22}% definite-theta-lifts-and-norms
\chapters@set{20261125T016500}{23}% hecke-operators-and-theta-lifts
\chapters@set{20261125T016700}{24}% divisibility-along-the-chain
\chapters@set{20261125T004000}{25}% denominators-and-sha
\chapters@set{20261125T004100}{26}% tate-twists-and-sha
\chapters@set{20261126T020430}{27}% the-minus-part-by-gauss-sums
\chapters@set{20261125T014100}{28}% congruence-modules
\chapters@set{20261125T014500}{29}% eisenstein-congruences-and-selmer-classes
\chapters@set{20261125T014700}{30}% pseudocharacters-and-the-length-bound
\chapters@set{20261125T011000}{31}% galois-representations-black-box
\chapters@set{20261126T030300}{32}% the-data
\chapters@set{20261126T020210}{33}% bessel-poincare-series-and-their-gram-operators
\chapters@set{20261126T030400}{34}% the-seesaw-identity
\chapters@set{20261126T030500}{35}% divisibility-of-the-composite-lift
\chapters@set{20261126T030600}{36}% theorem-a-short-forms
\chapters@set{20261126T030700}{37}% the-determinant-of-a-chain
\chapters@set{20261126T030800}{38}% the-cost-of-one-lift
\chapters@set{20261126T030900}{39}% denominators-of-projections-to-lifts
\chapters@set{20261126T031000}{40}% twists-with-few-congruences
\chapters@set{20261126T031100}{41}% theorem-c-and-the-main-inequality
\providecommand\chapnumber[1]{%
  \@ifundefined{chapters@#1}%
    {\PackageError{chapters}{Sheet #1 is not listed in the index}%
      {List it in 20261125T012000--grh-workbook__guide.org\MessageBreak
       and run scripts/make-chapters.py.}}%
    {\edef\thechap{\@nameuse{chapters@#1}}}}
\makeatother
\chapnumber{<sheet ID>}
\makeatletter
\providecommand\chapters@set[2]{\@namedef{chapters@#1}{#2}}
\chapters@set{20261125T001000}{1}% zeros-to-ramp-sums
\chapters@set{20261125T001100}{2}% inert-primes-and-estimate-b
\chapters@set{20261125T005000}{3}% galois-fourier-and-averaging
\chapters@set{20261125T005500}{4}% columns-from-a-sphere
\chapters@set{20261125T000000}{5}% slope-method
\chapters@set{20261126T030100}{6}% a-twist-for-each-prime
\chapters@set{20261125T003000}{7}% the-chain-and-the-square-root
\chapters@set{20261125T010000}{8}% rank-lemma-and-zero-estimate
\chapters@set{20261125T010500}{9}% poincare-series-and-near-orthogonality
\chapters@set{20261124T223000}{10}% elliptic-curves-for-the-proof
\chapters@set{20261125T009500}{11}% fock-model-and-theta-intertwiners
\chapters@set{20261124T223100}{12}% complex-multiplication-and-inert-reduction
\chapters@set{20261125T008000}{13}% one-cm-determinant
\chapters@set{20261124T214500}{14}% supersingular-theta-divisibility
\chapters@set{20261124T214600}{15}% the-isotropic-plane
\chapters@set{20261125T013700}{16}% pairings-on-a-formal-group-of-height-two
\chapters@set{20261125T013500}{17}% hecke-nondegeneracy-on-the-tree
\chapters@set{20261126T030200}{18}% the-bound-is-attained
\chapters@set{20261125T009000}{19}% kronecker-limit-and-chowla-selberg
\chapters@set{20261125T002000}{20}% definite-theta-lifts
\chapters@set{20261125T015500}{21}% siegel-weil-in-the-smallest-case
\chapters@set{20261125T002100}{22}% definite-theta-lifts-and-norms
\chapters@set{20261125T016500}{23}% hecke-operators-and-theta-lifts
\chapters@set{20261125T016700}{24}% divisibility-along-the-chain
\chapters@set{20261125T004000}{25}% denominators-and-sha
\chapters@set{20261125T004100}{26}% tate-twists-and-sha
\chapters@set{20261126T020430}{27}% the-minus-part-by-gauss-sums
\chapters@set{20261125T014100}{28}% congruence-modules
\chapters@set{20261125T014500}{29}% eisenstein-congruences-and-selmer-classes
\chapters@set{20261125T014700}{30}% pseudocharacters-and-the-length-bound
\chapters@set{20261125T011000}{31}% galois-representations-black-box
\chapters@set{20261126T030300}{32}% the-data
\chapters@set{20261126T020210}{33}% bessel-poincare-series-and-their-gram-operators
\chapters@set{20261126T030400}{34}% the-seesaw-identity
\chapters@set{20261126T030500}{35}% divisibility-of-the-composite-lift
\chapters@set{20261126T030600}{36}% theorem-a-short-forms
\chapters@set{20261126T030700}{37}% the-determinant-of-a-chain
\chapters@set{20261126T030800}{38}% the-cost-of-one-lift
\chapters@set{20261126T030900}{39}% denominators-of-projections-to-lifts
\chapters@set{20261126T031000}{40}% twists-with-few-congruences
\chapters@set{20261126T031100}{41}% theorem-c-and-the-main-inequality
\providecommand\chapnumber[1]{%
  \@ifundefined{chapters@#1}%
    {\PackageError{chapters}{Sheet #1 is not listed in the index}%
      {List it in 20261125T012000--grh-workbook__guide.org\MessageBreak
       and run scripts/make-chapters.py.}}%
    {\edef\thechap{\@nameuse{chapters@#1}}}}
\makeatother

\ifcsname chapters@20261126T020140\endcsname
  \chapnumber{20261126T020140}
\else
  \renewcommand{\thechap}{?}
\fi
\renewcommand{\thesection}{\thechap.\arabic{section}}
\theoremstyle{definition}
\newtheorem{problem}{Problem}[section]
\theoremstyle{plain}
\newtheorem{theorem}{Theorem}[section]
\theoremstyle{remark}
\newtheorem*{leanremark}{In the Lean formalization}
\newcommand{\Z}{\mathbb Z}
\newcommand{\GL}{\mathrm{GL}}
\newcommand{\LP}{\mathrm{LP}}
\newcommand{\rk}{\operatorname{rk}}
\newcommand{\len}{\operatorname{length}}
\newcommand{\vo}{\mathfrak o}
\newcommand{\Vc}{\mathcal V}
\newcommand{\tl}[1]{\texttt{#1}}
\setlength{\emergencystretch}{3em}

\title{The colength of a Hecke image}
\author{}
\date{November 26, 2026}

\begin{document}
\maketitle
\refstepcounter{chap}\label{chap}

\paragraph{Prerequisites.}
Modules over a discrete valuation ring and the elementary divisor theorem \cite[Ch.~III, \S7]{La}.  Chapter~\ref{hn-chap}, Theorem~\ref{hn-thm:vertex}; Chapter~\ref{ea-chap}: edges, the lattices \(H(L)\), and the condition \(\LP\) of Section~\ref{ea-sec:LP}.  Theorem~\ref{di-thm:colength} of Chapter~\ref{di-chap}, which states the result of this chapter in the form of \cite[Thm.~12.18]{Adr}.

Chapter~\ref{ea-chap} asked, at one edge, for vectors whose pairings are nondegenerate up to \(p^D\) on all but a fraction \(\kappa\) of the dimension.  What does that buy for the image of a product of Hecke operators, acting on forms over a whole class set and at several Frobenius orbits at once?  The answer is a bound for the colength of the image that is linear in the depth and the denominator, plus a term in \(\sum\kappa_v\).  The counting behind it is elementary, and we prove it in a model; the covering lemma that feeds it is stated.

Throughout, \(\vo\) is a complete discrete valuation ring with uniformizer \(\varpi\) and valuation \(v\), and \(p\in\vo\) is a nonzero nonunit; in the application \(v(p)=2\).  Modules are finitely generated, and \(\rk\) is the rank of a free \(\vo\)-module.

\section{Covering up to a few vectors}\label{sec:cover}

Let \(M\) be a free \(\vo\)-module and \(P\subseteq M\) a submodule.  For an integer \(E\ge0\) and a real \(g\ge0\) we say that \emph{\(p^EM\subseteq P\) up to \(g\) vectors} if there is a set \(C\subseteq M\) with at most \(g\) elements and \(p^EM\subseteq P+\vo C\).  If \(P\) has full rank, write \(P=\bigoplus_i\varpi^{\varepsilon_i}\vo h_i\) for a basis \((h_i)\) of \(M\), with exponents \(\varepsilon_i\).

\begin{problem}\label{pr:cover-smith}
Let \(P\) have full rank.  Then \(p^EM\subseteq P\) up to \(g\) vectors exactly when \(\#\{i:\varepsilon_i>Ev(p)\}\le g\).
\end{problem}

\begin{problem}\label{pr:length-cut}
Let \(P\) have full rank and \(\tau\ge0\).  Then \(P\cap p^\tau M=\bigoplus_i\varpi^{\max(\varepsilon_i,\tau v(p))}\vo h_i\), so \(\len(M/(P\cap p^\tau M))=\sum_i\max(\varepsilon_i,\tau v(p))\).
\end{problem}

\begin{problem}\label{pr:length-from-cover}
Let \(P\) have full rank, let \(\tau\ge0\) and \(S\ge1\) be integers, and suppose that \(p^{\tau+\beta S}M\subseteq P\) up to \(g_\beta\) vectors for every \(\beta\ge0\).  Then
\[
 \len\bigl(M/(P\cap p^\tau M)\bigr)\le v(p)\Bigl(\tau\rk M+S\sum_{\beta\ge0}g_\beta\Bigr).
\]
\end{problem}

\section{Elementary symmetric functions}\label{sec:sym}

For real numbers \(\kappa_v\ge0\), \(v\in\Vc\), let \(e_j(\kappa)=\sum_{U\subseteq\Vc,\,\lvert U\rvert=j}\prod_{v\in U}\kappa_v\).

\begin{problem}\label{pr:maclaurin}
\(e_j(\kappa)\le(\sum_v\kappa_v)^j/j!\) for every \(j\ge0\).
\end{problem}

\begin{problem}\label{pr:exp}
For \(0\le\Sigma\le1\), \(\sum_{\beta\ge0}\Sigma^{\beta+1}/(\beta+1)!=e^\Sigma-1\le2\Sigma\).  The inequality \(e^\Sigma-1\le2\Sigma\) fails for \(\Sigma=1.26\).
\end{problem}

\begin{problem}\label{pr:combine}
Let \(m\ge1\) and \(\Sigma=m\sum_v\kappa_v\le1\).  In Problem~\ref{pr:length-from-cover} suppose \(g_\beta\le\rk M\,m^{\beta+1}e_{\beta+1}(\kappa)\) for every \(\beta\).  Then
\[
 \len\bigl(M/(P\cap p^\tau M)\bigr)\le v(p)\rk M\,(\tau+2S\Sigma).
\]
\end{problem}

\section{Hecke operators on a tensor product}\label{sec:tensor}

Here is the mechanism in a model.  Let \(\Vc\) be a finite set, \(M_v\) free \(\vo\)-modules for \(v\in\Vc\), and \(M=\bigotimes_{v\in\Vc}M_v\).  Let \(Y_v\) be an endomorphism of \(M_v\), and \(y_v\) the endomorphism of \(M\) acting by \(Y_v\) on the factor \(M_v\) and by the identity on the others.  Suppose that for each \(v\) there are an integer \(E_v\ge0\) and a real \(\kappa_v\ge0\) with
\[
 p^{E_v}M_v\subseteq Y_vM_v\quad\text{up to }\kappa_v\rk M_v\text{ vectors}.
\]

\begin{problem}\label{pr:powers}
For \(m\ge1\), \(p^{mE_v}M_v\subseteq Y_v^mM_v\) up to \(m\kappa_v\rk M_v\) vectors.
\end{problem}

\begin{problem}\label{pr:expand}
For submodules \(A_v,B_v\subseteq M_v\), the image of \(\bigotimes_v(A_v+B_v)\) in \(M\) is the sum, over \(U\subseteq\Vc\), of the images of \(\bigotimes_{v\notin U}A_v\otimes\bigotimes_{v\in U}B_v\).
\end{problem}

Fix \(m\ge1\), integers \(0\le S_v\le S\), and put
\[
 N=\sum_{U\subseteq\Vc}\Bigl(\prod_{v\notin U}y_v^m\Bigr)\Bigl(p^{\sum_{v\in U}S_v}M\Bigr),\qquad \tau=m\sum_{v}E_v .
\]
So \(N\) is spanned by the images of products of Hecke operators, where each orbit not used is replaced by a scalar \(p^{S_v}\).

\begin{problem}\label{pr:tensor-count}
For every \(\beta\ge0\), \(p^{\tau+\beta S}M\subseteq N\) up to \(\rk M\,m^{\beta+1}e_{\beta+1}(\kappa)\) vectors.  Also \(p^{\sum_vS_v}M\subseteq N\).
\end{problem}

\begin{problem}\label{pr:sums}
If \(M=\bigoplus_cM^{(c)}\) is a direct sum of such tensor products, with the same \(\Vc\), \(E_v\), \(\kappa_v\), \(m\) and \(S_v\), and \(y_v\) acts on each summand as above, then Problem~\ref{pr:tensor-count} holds for \(M\).
\end{problem}

\begin{problem}\label{pr:example}
Let \(\vo=\Z_p\), \(\Vc=\{1,2\}\), \(M_v=\Z_p^4\), \(Y_v=\operatorname{diag}(1,1,1,p^5)\), \(E_v=0\), \(\kappa_v=1/4\), \(m=1\) and \(S_v=2\).  Then the exponents of \(N\) in \(M\) are \(0\) (nine times), \(2\) (six times) and \(4\) (once), so \(\len(M/N)=16\).  The bounds of Problem~\ref{pr:tensor-count} are \(8\) for \(\beta=0\) and \(1\) for \(\beta=1\); Problem~\ref{pr:length-from-cover} gives \(\len(M/N)\le18\), and Problem~\ref{pr:combine}, with \(\Sigma=1/2\), gives \(32\).
\end{problem}

The count \(\rk M\,m^{\beta+1}e_{\beta+1}(\kappa)\) is the number of basis vectors that are bad at \(\beta+1\) orbits at once.  A vector bad at only \(\beta\) orbits is rescued by the scalars \(p^{S_v}\) there, at the cost \(\beta S\) in depth.

\section{Forms on a class set}\label{sec:forms}

At one prime \(p\), let \(\Vc\) be a finite set of Frobenius orbits.  For each \(v\in\Vc\) take the data of Chapter~\ref{ea-chap}: \(W_v\), \(K_v\), a rank \(k\), a type \(r\) of edges, a representation \(V_v\) with lattices \(H_v(L)\) for \(W_v\)-lattices \(L\subset K_v^k\), and the normalizing exponent \(A_v\), with \(p^{A_v}H_v(L'')\subseteq H_v(L)\) for every edge (Problem~\ref{ea-pr:inclusions}).  A \emph{tuple} is a family \(L=(L_v)_v\) of lattices, and \(H(L)=\bigotimes_vH_v(L_v)\).  Let \(\Gamma\) be a group acting on tuples through homomorphisms \(\Gamma\to\GL_k(K_v)\), and \(\chi\colon\Gamma\to\vo^\times\) a character.  A \emph{form} is a function \(F\) on tuples with \(F(L)\in H(L)\) and \(F(\gamma L)=\chi(\gamma)\,\gamma F(L)\); write \(M\) for the \(\vo\)-module of forms.  Assume that \(\Gamma\) has finitely many orbits on tuples, with representatives \(L^{(1)},\dots,L^{(h)}\), and that the stabilizer in \(\Gamma\) of every tuple is trivial.

\begin{problem}\label{pr:delta}
\(F\mapsto(F(L^{(c)}))_c\) is an isomorphism \(M\cong\bigoplus_cH(L^{(c)})\).  Hence \(M\) is free of rank \(h\prod_v\dim V_v\).
\end{problem}

The normalized Hecke operator at \(v\) is
\[
 (y_vF)(L)=p^{A_v}\sum_{L''}F(L_{v\to L''}),
\]
the sum over the \(r\)-edges \((L''\subset L_v)\), where \(L_{v\to L''}\) is \(L\) with \(L_v\) replaced by \(L''\).  For \(\Vc\) a single orbit, \(k=2\), \(r=1\) and \(\nu=1\), this is the operator \(Y=p^kT\) of Chapter~\ref{hn-chap}, with \(k\) there the degree.

\begin{problem}\label{pr:hecke}
\(y_v\) maps \(M\) into itself, and \(y_uy_v=y_vy_u\) for all \(u,v\in\Vc\).
\end{problem}

In this setting the forms are not a tensor product over the orbits on which the \(y_v\) act factorwise, since \(y_v\) moves a tuple to its neighbours.  The formalization proves the covering of Section~\ref{sec:tensor} directly.

\begin{theorem}[colength of the Hecke image]\label{thm:colength}
In the setting above, let \(\Vc'\subseteq\Vc\), let \(m\ge1\), \(D\ge0\), \(c_*\ge0\), integers \(t_v\ge0\) and \(0\le S_v\le S\), and reals \(0\le\kappa_v\le1\), and suppose that \(\LP(r;t_v,D,\kappa_v,c_*)\) holds at every \(v\in\Vc'\).  Put
\[
 N=\sum_{U\subseteq\Vc'}\Bigl(\prod_{v\in\Vc'\setminus U}y_v^m\Bigr)\Bigl(p^{\sum_{v\in U}S_v}M\Bigr),\qquad \tau=m\sum_{v\in\Vc'}(t_v+D+1),\qquad \Sigma=m\sum_{v\in\Vc'}\kappa_v .
\]
\begin{itemize}
\item[(a)] For every \(\beta\ge0\), \(p^{\tau+\beta S}M\subseteq N\) up to \(\lfloor\rk M\,m^{\beta+1}e_{\beta+1}(\kappa)\rfloor\) vectors, and \(p^{\sum_vS_v}M\subseteq N\).
\item[(b)] If \(\Sigma\le1/2\), then \(\len(M/(N\cap p^\tau M))\le v(p)\rk M\,(\tau+2S\Sigma)\).  If \(M'\subseteq M\) is a saturated submodule, then \(N\cap p^\tau M\cap M'\subseteq p^\tau M'\), and \(\len(M'/(N\cap p^\tau M\cap M'))\) satisfies the same bound.
\end{itemize}
\end{theorem}

\noindent\emph{Status}: proved in the Lean formalization \cite{Lean}.  The proof of~(a) first covers \(p^{D+1}\) times a depth-\(t\) lattice of forms by Hecke images at the orbits of any \(U\), up to \(\prod_{v\in U}\kappa_v\rk M\) vectors.  For this it uses the isomorphism of Problem~\ref{pr:delta}, the isolated vectors of the condition \(\LP\), which make \(y_v\) of a form concentrated at a source almost concentrated at one target, and the witnesses and tests of~(iii) there.  It then counts as in Section~\ref{sec:tensor}.  We have not followed the proof of this covering lemma.  Part~(b) follows from~(a) by the problems below.  The corresponding statement in \cite{Adr} is \cite[Thm.~12.18]{Adr}, recalled in Theorem~\ref{di-thm:colength}.

\begin{problem}\label{pr:a-to-b}
Part~(a) of Theorem~\ref{thm:colength} implies the length bound in part~(b).
\end{problem}

\begin{problem}\label{pr:saturated}
Let \(M'\subseteq M\) be saturated, that is, \(M\cap(M'\otimes\mathrm{Frac}\,\vo)=M'\).  Then \(p^\tau M\cap M'=p^\tau M'\), and \(M'/(Q\cap M')\) injects into \(M/Q\) for every submodule \(Q\subseteq M\).
\end{problem}

The formalization uses two integral structures, \(M\subseteq M^{\rm na}\), on the forms at a prime; Theorem~\ref{thm:colength} is proved for \(M^{\rm na}\) and transferred.

\begin{problem}\label{pr:two-structures}
Let \(p^{s}M^{\rm na}\subseteq M\subseteq M^{\rm na}\), both stable under the \(y_v\), and write \(N(\cdot)\) for the module \(N\) built from each.  If \(p^EM^{\rm na}\subseteq N(M^{\rm na})\) up to \(g\) vectors, then \(p^{E+s}M\subseteq N(M)\) up to \(g\) vectors.  Hence Theorem~\ref{thm:colength} holds for \(M\) with \(\tau\) replaced by \(\tau+s\).
\end{problem}

With \(s=e_p\), \(D=D_p\), and the real depths \(t_v\) of Chapter~\ref{di-chap} replaced by \(\lceil t_v\rceil\), this is the exponent \(\tau(p)=e_p+m\sum_v(\lceil t_v\rceil+D_p+1)\) of Chapter~\ref{di-chap}, and \(2S\Sigma\), with \(S=\max_vS_v\), is the term \(\mathrm{io}(p)=2\Sigma(p)\max_vS_v\) there.  In the formalization the local lattice of the source at a served place \(w\mid p\) is \(N\cap p^{\lceil\tau(p)\rceil}M\cap V^{\rm ram}\), and Theorem~\ref{thm:colength}(b), with \(M'=M\cap V^{\rm ram}\), is Theorem~\ref{di-thm:colength}.

\begin{problem}\label{pr:tree}
Let \(P\subseteq M\) have full rank and \(p^EM\subseteq P\).  Then \(\len(M/P)\le v(p)E\rk M\).  Over \(W\), where \(v(p)=1\), with \(M=H(L)\) and \(P=S(L)\) on the tree of Chapter~\ref{hn-chap} with \(Q>d\), where \(d\) is the degree, Theorem~\ref{hn-thm:vertex} gives \(E=\lfloor\log_pd\rfloor\); for \(p=d=5\) and \(Q=25\) the bound is \(6\), and the true loss is \(4\).
\end{problem}

\begin{problem}\label{pr:share-check}
Let \(\rk M=8\), \(v(p)=2\), \(m=2\), one served orbit with \(\kappa=1/8\) and \(S=3\).  Then \(\Sigma=1/4\), \(2S\Sigma=3/2\), and the bound of Theorem~\ref{thm:colength}(b) is \(16\tau+24\), as in Problem~\ref{di-pr:share}.
\end{problem}

The colength is paid once per prime: it is the price \(\tau(p)+\mathrm{io}(p)\) of recording, as a change of lattice, the gain that the chain of lifts makes on \(N\).  The gain itself is a separate statement about the lifts.

\begin{leanremark}
\raggedright In \cite{Lean} ``\(p^EM\subseteq P\) up to \(g\) vectors'' is \tl{AllButLE}.  Theorem~\ref{thm:colength} is \tl{SelectionLemmas.Conclusion1}, proved as \tl{LatticeSystems.theta\_\allowbreak{}kernel\_\allowbreak{}expansion} from the covering lemma \tl{RestJ}, itself built from \tl{JDelta} and \tl{JRank} (Problem~\ref{pr:delta}), \tl{JSmith} (the condition \(\LP\) at one orbit), \tl{JStep1} (isolation) and \tl{JLocal} (the product over orbits).  The module \(N\) is \tl{prodIdeal}, and \(y_v\) is \tl{GAtP.yv}.  Problem~\ref{pr:combine} is \tl{floor\_\allowbreak{}scaled\_\allowbreak{}esymmOn\_\allowbreak{}le\_\allowbreak{}sum\_\allowbreak{}pow} with the summation in \tl{core}; Problem~\ref{pr:two-structures} is \tl{core} and \tl{coreT}; the version at a place is \tl{FormalGroupSystem.LemmaVVolumeGeneral}, proved by \tl{stmt6G}, and it gives the lattice clause of \tl{LinkG}.
\end{leanremark}

\section*{Hints}

\subsection*{\S\ref{sec:cover}}
\begin{description}
\item[\ref{pr:cover-smith}] For one direction take \(C\) among the \(h_i\); for the other, count the generators of \((p^EM+P)/P\).
\item[\ref{pr:length-cut}] \(p^\tau M\) is diagonal in every basis of \(M\).
\item[\ref{pr:length-from-cover}] Problems~\ref{pr:cover-smith} and~\ref{pr:length-cut}; count each exponent by the levels \(\tau+\beta S\) it exceeds.
\end{description}

\subsection*{\S\ref{sec:sym}}
\begin{description}
\item[\ref{pr:maclaurin}] Expand \((\sum_v\kappa_v)^j\).
\item[\ref{pr:exp}] Convexity of \(e^\Sigma-1-2\Sigma\).
\item[\ref{pr:combine}] Problems~\ref{pr:maclaurin} and~\ref{pr:exp}.
\end{description}

\subsection*{\S\ref{sec:tensor}}
\begin{description}
\item[\ref{pr:powers}] Induction on \(m\), keeping the vectors \(Y_v^iC\).
\item[\ref{pr:tensor-count}] Problems~\ref{pr:powers} and~\ref{pr:expand} with \(A_v=Y_v^mM_v\); treat the terms by the size of \(U\).
\item[\ref{pr:example}] Work in the basis \(e_i\otimes e_j\); each of the four generators of \(N\) is diagonal.
\end{description}

\subsection*{\S\ref{sec:forms}}
\begin{description}
\item[\ref{pr:delta}] Trivial stabilizers make the extension from the representatives well defined.
\item[\ref{pr:hecke}] \(\gamma\) maps the edges into \(L_v\) to the edges into \(\gamma L_v\); for integrality use Problem~\ref{ea-pr:inclusions}.
\item[\ref{pr:a-to-b}] Problems~\ref{pr:length-from-cover} and~\ref{pr:combine}; if \(S=0\), then \(N=M\).
\item[\ref{pr:two-structures}] \(p^sN(M^{\rm na})\subseteq N(M)\).
\item[\ref{pr:tree}] Problems~\ref{pr:cover-smith} and~\ref{pr:length-cut}, and Problem~\ref{di-pr:vertex}.
\end{description}

\begin{thebibliography}{Lean}
\bibitem[Adr]{Adr}  J.~Adriazola, \emph{The Alp\"oge proof of the Riemann hypothesis}, manuscript, revised 24 November 2026, \url{https://github.com/jimjamzola/RH-Alpoge-Reconstruction}.
\bibitem[La]{La}  S.~Lang, \emph{Algebra}, revised 3rd ed., Graduate Texts in Mathematics~211, Springer, 2002.
\bibitem[Lean]{Lean} The Lean community, \emph{riemann-hypothesis}, a Lean formalization of Alp\"oge's argument, \url{https://github.com/leanprover-community/riemann-hypothesis}.
\end{thebibliography}
\end{document}
